\documentclass{article}

\usepackage{amsmath, amsthm, amssymb, amsfonts}
\usepackage{thmtools}
\usepackage{graphicx}
\usepackage{setspace}
\usepackage{geometry}
\usepackage{float}
\usepackage[backend=biber,style=ieee,sorting=ynt]{biblatex}
\usepackage[hidelinks]{hyperref}
\usepackage[utf8]{inputenc}
\usepackage[spanish]{babel}
\usepackage{framed}
\usepackage[dvipsnames]{xcolor}
\usepackage{tcolorbox}
\usepackage{tikz}
\usetikzlibrary{positioning}
\usetikzlibrary{arrows.meta}
\usepackage{caption}
\usepackage{longtable}
\usepackage{pdflscape}
\usepackage{svg}
\usepackage{subcaption}
\usepackage{multirow}
\usepackage{array}
\usepackage{listings}
\usepackage{cancel}
\usepackage{fancyhdr}

\addbibresource{referencias.bib}

\colorlet{LightGray}{White!90!Periwinkle}
\colorlet{LightOrange}{Orange!15}
\colorlet{LightGreen}{Green!15}

\newcommand{\HRule}[1]{\rule{\linewidth}{#1}}

\hypersetup{
    pdftitle={Propuesta de Proyecto: Agente Inteligente para BreakApp},
    pdfauthor={Equipo BreakApp},
    pdfsubject={Propuesta de Proyecto de Inteligencia Artificial},
    pdfkeywords={Inteligencia Artificial, Sistema de Recomendación, Agente Racional, PEAS, BreakApp},
    pdfcreator={LaTeX with hyperref},
    pdfproducer={pdflatex}
}

\setstretch{1.2}
\geometry{
    textheight=9in,
    textwidth=5.5in,
    top=1in,
    headheight=12pt,
    headsep=25pt,
    footskip=30pt
}

\pagestyle{fancy}
\fancyhf{}
\renewcommand{\headrulewidth}{0pt}
\fancyfoot[C]{\footnotesize \thepage \quad \textbar{} \quad Este trabajo está bajo una licencia CC BY-SA 4.0.\\ Más info: \url{https://creativecommons.org/licenses/by/4.0/}}

\begin{document}

% ------------------------------------------------------------------------------
% Cover Page and ToC
% ------------------------------------------------------------------------------

\title{ \normalsize \textsc{}
	\\ [2.0cm]
	\HRule{1.5pt} \\
	\LARGE \textbf{\uppercase{BreakApp-AI}
		\HRule{2.0pt} \\ [0.6cm] \LARGE{Agente racional para la recomendación contextual y optimización predictiva de pedidos en cafeterías universitarias} \vspace*{10\baselineskip}}
}
\date{}
\author{\textbf{Andrés Rubiano} \\
	 Inteligencia Artificial - 2026-2S \\
	 Universidad de Bogotá Jorge Tadeo Lozano}

\maketitle
\thispagestyle{empty}
\newpage

\tableofcontents
\thispagestyle{empty}
\newpage
\setcounter{page}{1}

\section{Descripción del problema}

En el contexto de la vida universitaria contemporánea, los estudiantes, docentes y personal administrativo disponen de recesos breves (\textit{breaks}) de entre 15 y 30 minutos entre jornadas de clase para descansar, alimentarse e interactuar socialmente. Sin embargo, la experiencia de adquisición de alimentos y bebidas dentro de los puntos de venta del campus (tales como \textit{Catering Gourmet}, \textit{Punto Verde} y \textit{Central Food}) se ve sistemáticamente degradada por una severa fricción operativa y cognitiva~\cite{ricci2022recommender}. 

Por un lado, los usuarios se enfrentan a menús estáticos, dispersos y con alta sobrecarga de información, lo que genera indecisión en el momento de la compra y pérdida de tiempo valioso en filas físicas o pantallas de selección. Asimismo, las preferencias individuales varían sustancialmente según factores multidimensionales: horarios de consumo (café y snacks matutinos entre las 8:00 AM y 8:30 AM vs. almuerzos balanceados entre las 12:30 PM y 1:30 PM), presupuestos semanales heterogéneos, hábitos nutricionales específicos (arquetipos fitness, veganos o estudiantes vespertinos) y restricciones dietéticas no negociables (intolerancia a la lactosa, exclusión de ingredientes específicos como cebolla o dietas bajas en azúcar)~\cite{adomavicius2011context}. 

Por otro lado, los restaurantes del campus operan con picos de demanda altamente concentrados y asimétricos, provocando congestión en cocina, agotamiento imprevisto de platos populares y mermas en productos perecederos. A pesar de contar con la plataforma móvil \textbf{BreakApp} para la visualización del menú y el trámite de compras, la aplicación opera bajo una presentación pasiva y determinista de productos, careciendo de un mecanismo automatizado e inteligente que anticipe las necesidades del usuario en función del tiempo, el contexto y el historial de consumo~\cite{burke2002hybrid}.

Ante la existencia de registros estructurados de órdenes, horarios, notas de personalización y catálogo enriquecido de productos en BreakApp, surge la oportunidad de formular e implementar un \textbf{Agente Racional} basado en Inteligencia Artificial~\cite{russell2020artificial}. Dicho agente percibirá el estado contextual del consumidor (identidad, historial, franja horaria, día de la semana y restricciones declaradas) y el estado de la oferta (disponibilidad, categoría complementaria y afinidades de maridaje), evaluará las alternativas mediante una función de utilidad multicriterio y actuará recomendando de forma proactiva el conjunto óptimo de productos y combos sugeridos, reduciendo la fricción en la toma de decisiones y equilibrando el flujo de pedidos hacia la cocina.

\section{Objetivos}

\subsection{General}

Diseñar e implementar un agente racional basado en inteligencia artificial que, a partir del análisis de datos históricos de consumo, atributos del catálogo y variables contextuales temporales en la aplicación BreakApp, genere recomendaciones predictivas y combinaciones óptimas de productos (combos/maridajes) con el fin de maximizar la utilidad percibida por el usuario y minimizar los tiempos de indecisión en los recesos de campus, sujeto a restricciones dietéticas y presupuestales.

\subsection{Específicos}

\begin{itemize}
    \item Formalizar el entorno y la estructura PEAS del agente (Rendimiento, Entorno, Actuadores y Sensores), delimitando con precisión el espacio de estados, acciones, supuestos operativos y restricciones duras de alérgenos y presupuesto a partir de los datos de \textit{BreakApp}~\cite{russell2020artificial}.
    \item Diseñar y calibrar una función de utilidad multicriterio que pondere matemáticamente la afinidad histórica por categorías y etiquetas (\textit{tags}), la congruencia temporal (hora del día y día de la semana), la penalización estricta por incumplimiento de restricciones dietéticas (\textit{notes}) y la afinidad de maridaje complementario (\textit{pairing})~\cite{adomavicius2011context}.
    \item Implementar un motor de recomendación contextual híbrido (Filtrado Basado en Contenido combinado con reglas de afinidad temporal y grafos de maridaje de productos) desacoplado de la interfaz móvil, capaz de procesar solicitudes de inferencia con tiempos de respuesta en milisegundos~\cite{burke2002hybrid}.
    \item Evaluar el desempeño del agente inteligente comparándolo contra políticas base de recomendación estática (como productos más vendidos o aleatorios), cuantificando la ganancia neta de utilidad, la precisión de acierto en el Top-$K$ y la adherencia estricta a restricciones alimentarias.
\end{itemize}

\section{Enlace vídeo}

Entrar con el correo institucional \texttt{utadeo.edu.co}:

\vspace{0.3cm}
\noindent \textbf{Enlace de sustentación:} \url{https://utadeoeduco0-my.sharepoint.com/}

\addcontentsline{toc}{section}{Referencias}
\printbibliography

\newpage
\thispagestyle{empty}
\vspace*{0.3\textwidth}

\begin{figure}[H]
  \centering
  \includesvg[width=\textwidth]{img/LogoUtadeo.svg}
\end{figure}

\end{document}
